\documentclass[12pt]{article}
\def\activityid{F06-U-v2}\usepackage[letterpaper,margin=.65in,top=.60in,bottom=.65in]{geometry}
\usepackage[T1]{fontenc}\usepackage[default]{sourcesanspro}
\usepackage{microtype,tikz,array,tabularx,fancyhdr,amsmath}\usepackage[hidelinks]{hyperref}
\usetikzlibrary{arrows.meta,calc}
\setlength{\parindent}{0pt}\setlength{\parskip}{7pt}\setlength{\headheight}{0pt}\setlength{\footskip}{26pt}\setlength{\emergencystretch}{2em}
\pagestyle{fancy}\fancyhf{}\renewcommand{\headrulewidth}{0pt}\renewcommand{\footrulewidth}{.4pt}
\fancyfoot[L]{\scriptsize Bellingham Math Circle / Week 6 / \activityid}\fancyfoot[R]{\scriptsize \thepage}
\definecolor{ink}{gray}{.12}\definecolor{soft}{gray}{.95}\color{ink}
\newcommand{\PageTitle}[2]{{\fontsize{10}{12}\selectfont\bfseries WEEK 6\quad CODE LAB\hfill #1\par}\vspace{3pt}{\fontsize{26}{29}\selectfont\bfseries #2\par}\vspace{2pt}}
\newcommand{\NameLine}{{\small Name: \rule{2.65in}{.4pt}\hfill Date: \rule{1.35in}{.4pt}\par}\vspace{4pt}}
\newcommand{\Task}[2]{\par\vspace{3pt}{\large\bfseries #1}\par #2\par}
\newcommand{\Subhead}[1]{\par\vspace{3pt}{\large\bfseries #1}\par}
\newcommand{\RuleBox}[1]{\par\begingroup\setlength{\fboxsep}{8pt}\colorbox{soft}{\parbox{\dimexpr\linewidth-16pt\relax}{#1}}\endgroup\par}
\newcommand{\WriteBox}[2]{\par\textbf{#1}\par\vspace{2pt}\begin{tikzpicture}\draw[black!40,line width=.5pt] (0,0) rectangle (\dimexpr\linewidth-.5pt\relax,#2);\end{tikzpicture}\par}
\newcommand{\StudentFooter}{\vfill{\small Build first. Explain by pointing, talking, or drawing.}\par}
\newcommand{\Code}[1]{\texttt{#1}}

\hypersetup{pdftitle={Week 6: grades-4-5}}
\begin{document}
\PageTitle{GRADES 4--5}{Build a guaranteed code breaker}\NameLine
\RuleBox{Two colors: R and B. Repeats allowed. A test earns one point for each color in the correct position. You hear \textbf{only the total}, not which positions match.}\Task{1. Play four places.}{Try a hidden code. Then build a method that always identifies it in at most \textbf{four scored tests}. You may name the answer without paying for a final test.}
\begin{center}\renewcommand{\arraystretch}{2}\begin{tabular}{c|p{2.2in}|p{.7in}|p{1.5in}}Test&My code&Score&Still possible\\\hline 1&&&\\2&&&\\3&&&\\4&&&\end{tabular}\end{center}
\Task{2. Explain a baseline method.}{Test \Code{RRRR}. Then separately change place 1, place 2, and place 3 to B, keeping every other place R. Explain how the fourth place is recovered.}
\WriteBox{What each score change tells me, and why the method covers every secret:}{1.05in}
\textbf{A real optimization question:} Four tests are enough. Could a cleverer, adaptive method always manage three? On the next page, play the role of a difficult secret maker.\StudentFooter
\newpage
\PageTitle{GRADES 4--5}{The secret maker fights back}\NameLine
\RuleBox{Two colors: R and B. Repeats allowed. A test earns one point for each color in the correct position. You hear \textbf{only the total}, not which positions match.}\Task{3. Make the first answer difficult.}{After \Code{RRRR}, answer 2. These are the six possible secrets. Group them by their score for each proposed second test.}
\begin{center}\large\Code{BBRR}\quad\Code{BRBR}\quad\Code{BRRB}\par\Code{RBBR}\quad\Code{RBRB}\quad\Code{RRBB}\end{center}
\begin{center}\renewcommand{\arraystretch}{2}\begin{tabular}{c|p{1.65in}|p{1.65in}|p{1in}}Second test&Score groups&Largest group&Its size\\\hline BRRR&&&\\BBRR&&&\\BBBR&&&\\BBBB&&&\end{tabular}\end{center}
\Task{4. Only one test remains.}{For \Code{BRRR}, let the answer leave three candidates. Can one last test separate those three? For \Code{BBRR}, let the answer leave four candidates. Can one last test separate those four?}
\WriteBox{Try final tests; explain why each difficult branch defeats them.}{.65in}
\textbf{Three-secret hint:} Use \Code{BBRR}, \Code{BRBR}, \Code{BRRB}. Which position is fixed? What is the only feature that changes? A test asks whether the moving B falls in one of its B positions. How many scores can that give?

\textbf{Four-secret hint:} All six original candidates contain exactly two Bs. How many scores can one test produce on them? Is that enough to separate four candidates?\StudentFooter
\newpage
\PageTitle{GRADES 4--5}{Turn the experiment into a proof}\NameLine
\Task{5. Cover every possible second test.}{Second tests have 0, 1, 2, 3, or 4 Bs. Changing position names does not change the problem. Swapping every R and B in a test replaces its score by \(4-\text{score}\). Use those facts to explain why the cases on page 2 cover everything.}
\WriteBox{Why one or three Bs leaves a hard three-secret branch:}{1.1in}
\WriteBox{Why two Bs leaves a hard four-secret branch:}{1.1in}
\Task{6. What about a different first test?}{Rename the two colors separately at selected positions, in both every secret and every test. Can you turn any first test into \Code{RRRR} without changing which places match?}
\WriteBox{My conclusion: the smallest guaranteed number of tests, and why:}{1in}
\RuleBox{A short method proves an \textbf{upper bound}. A difficult branch that defeats every shorter method proves a \textbf{lower bound}. When they meet, we know the optimum.}
\textbf{Go further:} Must five places need five tests? The extra page contains a surprise.\StudentFooter
\end{document}
